\documentclass[11pt,a4paper]{article}

\usepackage[margin=1in]{geometry}
\usepackage[T1]{fontenc}
\usepackage[utf8]{inputenc}
\usepackage{lmodern}
\usepackage{xcolor}
\usepackage{hyperref}
\usepackage{longtable}
\usepackage{array}
\usepackage{enumitem}
\usepackage{listings}
\usepackage{titlesec}

\hypersetup{
  colorlinks=true,
  linkcolor=blue,
  urlcolor=blue,
  citecolor=blue
}

\definecolor{codebg}{RGB}{247,247,247}
\definecolor{codeborder}{RGB}{220,220,220}
\definecolor{keyword}{RGB}{0,88,90}
\definecolor{comment}{RGB}{105,105,105}

\lstdefinestyle{jsstyle}{
  backgroundcolor=\color{codebg},
  basicstyle=\ttfamily\small,
  breaklines=true,
  frame=single,
  rulecolor=\color{codeborder},
  keywordstyle=\color{keyword}\bfseries,
  commentstyle=\color{comment},
  showstringspaces=false,
  tabsize=2
}

\lstset{style=jsstyle}
\setlist[itemize]{topsep=3pt,itemsep=2pt}
\setlist[enumerate]{topsep=3pt,itemsep=2pt}

\title{\textbf{RentPi HackSpark Project Logic Documentation}}
\author{Team Documentation}
\date{\today}

\begin{document}
\maketitle
\tableofcontents
\newpage

\section{Project Overview}

This project implements a microservice-based rental platform called \textbf{RentPi}. The original challenge asks teams to build a platform around a central historical data API. The system contains a frontend, an API gateway, several backend services, and MongoDB for local persistence.

\subsection{Main Goal}

The goal is to provide product browsing, authentication, rental analytics, product availability, seasonal recommendations, and an AI assistant that answers questions using real data.

\subsection{High-Level Architecture}

\begin{center}
\begin{tabular}{|p{3.2cm}|p{2cm}|p{8cm}|}
\hline
\textbf{Service} & \textbf{Port} & \textbf{Responsibility} \\
\hline
Frontend & 3000 & React interface for users, products, availability, analytics, and chat. \\
\hline
API Gateway & 8000 & Single public backend entry point. Proxies requests to internal services. \\
\hline
User Service & 8001 & Registration, login, JWT authentication, discount tier lookup. \\
\hline
Rental Service & 8002 & Product proxy, availability, busiest dates, top categories, free streak, merged feed. \\
\hline
Analytics Service & 8003 & Peak window, surge days, and seasonal recommendations. \\
\hline
Agentic Service & 8004 & AI chatbot with data grounding and MongoDB-backed chat history. \\
\hline
MongoDB & 27017 & Stores users, chat sessions, and chat messages. \\
\hline
\end{tabular}
\end{center}

\subsection{Central API}

The central API is read-only and provides historical RentPi data:

\begin{itemize}
  \item Categories
  \item Users and security scores
  \item Products
  \item Rentals
  \item Rental statistics by date or category
\end{itemize}

Every service that calls the central API attaches:

\begin{lstlisting}
Authorization: Bearer CENTRAL_API_TOKEN
\end{lstlisting}

\section{Request Flow}

All browser requests go through the API gateway:

\begin{lstlisting}
Browser -> frontend:3000 -> api-gateway:8000 -> target service
\end{lstlisting}

The frontend never calls the central API directly. This protects the team token and keeps backend behavior centralized.

\section{Environment Variables}

Important environment variables:

\begin{longtable}{|p{4cm}|p{10cm}|}
\hline
\textbf{Variable} & \textbf{Purpose} \\
\hline
\texttt{CENTRAL\_API\_URL} & Base URL of the central data API. \\
\hline
\texttt{CENTRAL\_API\_TOKEN} & Team token used for central API authorization. \\
\hline
\texttt{MONGO\_URI} & MongoDB connection string used by user-service and agentic-service. \\
\hline
\texttt{JWT\_SECRET} & Secret used to sign and verify JWT tokens. \\
\hline
\texttt{GEMINI\_API\_KEY} & API key for Gemini model used by the chat assistant. \\
\hline
\texttt{GEMINI\_MODEL} & Gemini model name, defaulting to \texttt{gemini-2.5-flash}. \\
\hline
\end{longtable}

\newpage
\section{Problem-by-Problem Logic}

\subsection{P1: Health Checks}

\subsubsection*{Goal}
Every backend service must expose a health endpoint at:

\begin{lstlisting}
GET /status
\end{lstlisting}

The gateway must aggregate all downstream service statuses.

\subsubsection*{Implemented Logic}

Each service returns:

\begin{lstlisting}
{
  "service": "service-name",
  "status": "OK"
}
\end{lstlisting}

The API gateway calls downstream services in parallel:

\begin{lstlisting}
Promise.all([
  fetch("http://user-service:8001/status"),
  fetch("http://rental-service:8002/status"),
  fetch("http://analytics-service:8003/status"),
  fetch("http://agentic-service:8004/status")
])
\end{lstlisting}

Final gateway response:

\begin{lstlisting}
{
  "service": "api-gateway",
  "status": "OK",
  "downstream": {
    "user-service": "OK",
    "rental-service": "OK",
    "analytics-service": "OK",
    "agentic-service": "OK"
  }
}
\end{lstlisting}

If a service cannot be reached, its status becomes:

\begin{lstlisting}
"UNREACHABLE"
\end{lstlisting}

\subsection{P2: User Authentication}

\subsubsection*{Goal}
Allow users to register, login, and view their authenticated profile.

\subsubsection*{Endpoints}

\begin{itemize}
  \item \texttt{POST /users/register}
  \item \texttt{POST /users/login}
  \item \texttt{GET /users/me}
\end{itemize}

\subsubsection*{Register Logic}

\begin{enumerate}
  \item Read \texttt{name}, \texttt{email}, and \texttt{password} from request body.
  \item Check if the email already exists in MongoDB.
  \item If duplicate, return \texttt{409 Conflict}.
  \item Hash password using \texttt{bcryptjs}.
  \item Save user to MongoDB.
  \item Generate JWT containing user id, name, and email.
  \item Return token to the client.
\end{enumerate}

\subsubsection*{Login Logic}

\begin{enumerate}
  \item Find user by email.
  \item If no user exists, return \texttt{401 Unauthorized}.
  \item Compare submitted password with hashed password.
  \item If password is wrong, return \texttt{401 Unauthorized}.
  \item Generate and return JWT.
\end{enumerate}

\subsubsection*{Protected Profile Logic}

The middleware reads:

\begin{lstlisting}
Authorization: Bearer <token>
\end{lstlisting}

Then it verifies the token using \texttt{JWT\_SECRET}. If valid, decoded user data is returned.

\subsection{P3: Product Proxy}

\subsubsection*{Goal}
Rental service proxies product requests to the central API while hiding the team token.

\subsubsection*{Endpoints}

\begin{itemize}
  \item \texttt{GET /rentals/products}
  \item \texttt{GET /rentals/products/:id}
\end{itemize}

\subsubsection*{Logic}

For list products:

\begin{lstlisting}
GET /rentals/products?category=TOOLS&page=2&limit=20
\end{lstlisting}

The service forwards all query parameters to:

\begin{lstlisting}
GET /api/data/products?category=TOOLS&page=2&limit=20
\end{lstlisting}

For product detail:

\begin{lstlisting}
GET /rentals/products/42
\end{lstlisting}

The service calls:

\begin{lstlisting}
GET /api/data/products/42
\end{lstlisting}

Errors from the central API are translated to useful API responses.

\subsection{P4: Docker Compose and Multistage Builds}

\subsubsection*{Goal}
Run the full stack using:

\begin{lstlisting}
docker compose up --build
\end{lstlisting}

\subsubsection*{Implemented Logic}

The project defines all services in \texttt{docker-compose.yml}. Each backend service has a Dockerfile with dependency installation and runtime stages.

For optimized service images:

\begin{itemize}
  \item Dependencies are installed with \texttt{npm ci --omit=dev}.
  \item Runtime images use lean Alpine with Node.js.
  \item \texttt{.dockerignore} excludes \texttt{node\_modules}, logs, markdown, and unnecessary folders.
  \item Runtime stages copy only required app files.
\end{itemize}

\subsection{P5: Paginated Product Listing with Category Filter}

\subsubsection*{Goal}
List products with pagination and optional category validation.

\subsubsection*{Endpoint}

\begin{lstlisting}
GET /rentals/products?category=TOOLS&page=2&limit=20
\end{lstlisting}

\subsubsection*{Logic}

\begin{enumerate}
  \item If category is provided, fetch valid categories from central API.
  \item Cache category list for one hour.
  \item If category is invalid, return \texttt{400} with all valid categories.
  \item Forward all query parameters to central products endpoint.
  \item Return the central API pagination envelope unchanged.
\end{enumerate}

\subsubsection*{Why Cache Categories?}

The central API has rate limits. Without caching, every product list request with category would make an extra central API call. Caching reduces unnecessary requests.

\subsection{P6: Loyalty Discount}

\subsubsection*{Goal}
Return a user's discount tier based on security score.

\subsubsection*{Endpoint}

\begin{lstlisting}
GET /users/:id/discount
\end{lstlisting}

\subsubsection*{Logic}

\begin{enumerate}
  \item Fetch user from central API.
  \item Read \texttt{securityScore}.
  \item Apply discount tier mapping.
  \item Return user id, score, and discount percent.
\end{enumerate}

\subsubsection*{Discount Mapping}

\begin{center}
\begin{tabular}{|c|c|}
\hline
\textbf{Security Score} & \textbf{Discount} \\
\hline
80--100 & 20\% \\
\hline
60--79 & 15\% \\
\hline
40--59 & 10\% \\
\hline
20--39 & 5\% \\
\hline
0--19 & 0\% \\
\hline
\end{tabular}
\end{center}

\subsection{P7: Product Availability}

\subsubsection*{Goal}
Check whether a product is available in a requested date range.

\subsubsection*{Endpoint}

\begin{lstlisting}
GET /rentals/products/:id/availability?from=2024-03-01&to=2024-03-14
\end{lstlisting}

\subsubsection*{Core Algorithm}

The main challenge is overlapping rental periods. The service solves this using interval merging.

\begin{enumerate}
  \item Fetch all rentals for product using pagination.
  \item Keep rentals that overlap the requested range.
  \item Convert each rental to interval \texttt{[start, end]}.
  \item Sort intervals by start date.
  \item Merge overlapping intervals.
  \item Scan merged intervals to compute free windows.
  \item If no busy period overlaps, product is available.
\end{enumerate}

\subsubsection*{Interval Merge Example}

Input:

\begin{lstlisting}
[2024-03-01, 2024-03-05]
[2024-03-04, 2024-03-08]
[2024-03-12, 2024-03-15]
\end{lstlisting}

Merged:

\begin{lstlisting}
[2024-03-01, 2024-03-08]
[2024-03-12, 2024-03-15]
\end{lstlisting}

\subsection{P8: Kth Busiest Date}

\subsubsection*{Goal}
Find the kth busiest rental date across a month range.

\subsubsection*{Endpoint}

\begin{lstlisting}
GET /rentals/kth-busiest-date?from=2024-01&to=2024-06&k=3
\end{lstlisting}

\subsubsection*{Validation}

\begin{itemize}
  \item \texttt{from} and \texttt{to} must be valid \texttt{YYYY-MM}.
  \item \texttt{k} must be a positive integer.
  \item \texttt{from} must not be after \texttt{to}.
  \item Range must be at most 12 months.
\end{itemize}

\subsubsection*{Algorithm}

Instead of sorting every day, the service uses a min-heap of size \texttt{k}.

\begin{enumerate}
  \item Fetch rental stats for each month.
  \item For each day-count pair:
  \begin{itemize}
    \item If heap size is less than \texttt{k}, insert the day.
    \item Else if current count is greater than heap minimum, remove minimum and insert current day.
  \end{itemize}
  \item After processing all days, heap root is the kth busiest day.
\end{enumerate}

\subsubsection*{Complexity}

\begin{lstlisting}
Time: O(N log K)
Space: O(K)
\end{lstlisting}

This is better than sorting all dates, which would be:

\begin{lstlisting}
O(N log N)
\end{lstlisting}

\subsection{P9: Renter Top Categories}

\subsubsection*{Goal}
Find top categories rented by a user.

\subsubsection*{Endpoint}

\begin{lstlisting}
GET /rentals/users/:id/top-categories?k=5
\end{lstlisting}

\subsubsection*{Logic}

\begin{enumerate}
  \item Fetch all rentals by renter id.
  \item Count rentals per product id.
  \item Deduplicate product ids.
  \item Fetch product details using batch endpoint, max 50 ids per call.
  \item Map product id to category.
  \item Count rentals per category.
  \item Return top \texttt{k} categories.
\end{enumerate}

\subsubsection*{Why Batch Product Fetching?}

The central API supports:

\begin{lstlisting}
GET /api/data/products/batch?ids=1,2,3
\end{lstlisting}

Using this avoids one request per product, reducing latency and avoiding rate-limit problems.

\subsection{P10: Longest Free Streak}

\subsubsection*{Goal}
Find the longest continuous period in a year when a product was not rented.

\subsubsection*{Endpoint}

\begin{lstlisting}
GET /rentals/products/:id/free-streak?year=2023
\end{lstlisting}

\subsubsection*{Logic}

\begin{enumerate}
  \item Validate year.
  \item Fetch all rentals for product.
  \item Keep rentals that overlap the target year.
  \item Clip rentals to year boundaries.
  \item Sort busy periods by start date.
  \item Merge overlapping or adjacent busy periods.
  \item Scan gaps between busy periods.
  \item Return the longest inclusive free streak.
\end{enumerate}

\subsubsection*{Important Boundary Rule}

Rental dates are treated as occupied days. Therefore:

\begin{itemize}
  \item A free period before a rental ends at one day before rental start.
  \item A free period after a rental begins one day after rental end.
\end{itemize}

Example:

\begin{lstlisting}
Busy: 2023-04-10 to 2023-04-12
Free before busy period ends at: 2023-04-09
Free after busy period starts at: 2023-04-13
\end{lstlisting}

\subsection{P11: Seven-Day Peak Window}

\subsubsection*{Goal}
Find the 7-day window with the highest total rentals in a month range.

\subsubsection*{Endpoint}

\begin{lstlisting}
GET /analytics/peak-window?from=2024-01&to=2024-06
\end{lstlisting}

\subsubsection*{Logic}

\begin{enumerate}
  \item Validate month range.
  \item Fetch rental date stats for every month.
  \item Create full calendar array from first day of \texttt{from} to last day of \texttt{to}.
  \item Fill missing dates with count \texttt{0}.
  \item Use sliding window of size 7.
  \item Track maximum sum and best start date.
\end{enumerate}

\subsubsection*{Sliding Window}

Initial sum:

\begin{lstlisting}
sum = day1 + day2 + ... + day7
\end{lstlisting}

Move one day:

\begin{lstlisting}
sum = sum - old_left_day + new_right_day
\end{lstlisting}

\subsubsection*{Complexity}

\begin{lstlisting}
Time: O(N)
Space: O(N)
\end{lstlisting}

\subsection{P12: Unified Feed}

\subsubsection*{Goal}
Merge rental histories for multiple products into one chronological feed.

\subsubsection*{Endpoint}

\begin{lstlisting}
GET /rentals/merged-feed?productIds=12,47,88&limit=30
\end{lstlisting}

\subsubsection*{Validation}

\begin{itemize}
  \item Product ids must be 1--10 comma-separated integers.
  \item Duplicates are deduplicated.
  \item Limit must be positive and at most 100.
\end{itemize}

\subsubsection*{Algorithm}

Each product rental history is already sorted by \texttt{rentalStart}. The service merges sorted lists using two pointers and divide-and-conquer merging.

\begin{enumerate}
  \item Fetch first \texttt{limit} rentals for each product id.
  \item Treat each product rental list as one sorted stream.
  \item Merge pairs of sorted streams.
  \item Keep only first \texttt{limit} records.
\end{enumerate}

\subsubsection*{Two-Pointer Merge}

\begin{lstlisting}
while both lists have records:
  compare rentalStart
  push earlier record
  advance that list pointer
\end{lstlisting}

\subsection{P13: Surge Days}

\subsubsection*{Goal}
For each day in a month, find the next future day with a strictly higher rental count.

\subsubsection*{Endpoint}

\begin{lstlisting}
GET /analytics/surge-days?month=2024-03
\end{lstlisting}

\subsubsection*{Logic}

\begin{enumerate}
  \item Validate month.
  \item Fetch rental stats for month.
  \item Build complete calendar for month.
  \item Fill missing dates with count \texttt{0}.
  \item Use monotonic stack to find next greater day.
\end{enumerate}

\subsubsection*{Monotonic Stack Explanation}

The stack stores days waiting for a future higher count. When current day count is greater than stack top, that current day is the answer for the popped day.

\begin{lstlisting}
for each day i:
  while stack not empty and count[i] > count[stack.top]:
    popped = stack.pop()
    answer[popped] = i
  stack.push(i)
\end{lstlisting}

\subsubsection*{Complexity}

\begin{lstlisting}
Time: O(N)
Space: O(N)
\end{lstlisting}

\subsection{P14: Seasonal Recommendations}

\subsubsection*{Goal}
Recommend products historically rented around the same date in the past two years.

\subsubsection*{Endpoint}

\begin{lstlisting}
GET /analytics/recommendations?date=2024-06-15&limit=10
\end{lstlisting}

\subsubsection*{Logic}

\begin{enumerate}
  \item Validate date and limit.
  \item For each of the past two years:
  \begin{itemize}
    \item Build 15-day seasonal window: 7 days before and 7 days after target date.
  \end{itemize}
  \item Fetch all rentals in those windows.
  \item Count frequency per product id.
  \item Use bucket sorting by frequency to select top products.
  \item Enrich top product ids using batch product endpoint.
  \item Return product id, name, category, and score.
\end{enumerate}

\subsubsection*{Why Use a Date Library?}

Date windows can cross month or year boundaries. For example:

\begin{lstlisting}
date = 2024-01-03
window = 2023-12-27 to 2024-01-10
\end{lstlisting}

The implementation uses \texttt{dayjs} to handle this safely.

\subsection{P15: RentPi Assistant}

\subsubsection*{Goal}
Build an AI assistant that answers RentPi questions using grounded platform data.

\subsubsection*{Endpoint}

\begin{lstlisting}
POST /chat
Body: { "sessionId": "abc123", "message": "Which category had the most rentals?" }
\end{lstlisting}

\subsubsection*{Main Rules}

\begin{itemize}
  \item Refuse unrelated questions before calling the LLM.
  \item Fetch real data for supported RentPi questions.
  \item Inject fetched data into model prompt.
  \item Never invent numbers when data is missing.
\end{itemize}

\subsubsection*{Topic Guard}

The service first checks whether the message contains RentPi-related keywords:

\begin{lstlisting}
rent, rental, product, category, price, availability,
discount, trend, recommendation, peak, surge
\end{lstlisting}

If the message is unrelated, it returns a fixed refusal and does not call Gemini.

\subsubsection*{Intent Classification}

The service classifies user messages into intents:

\begin{itemize}
  \item Most rented category
  \item Product availability
  \item Recommendations
  \item Peak window
  \item Surge days
  \item Discount tier
  \item Product detail
  \item General RentPi help
\end{itemize}

\subsubsection*{Grounding}

For each intent, the service fetches data from either:

\begin{itemize}
  \item Central API
  \item Rental service
  \item Analytics service
  \item User service
\end{itemize}

Then it sends the data to Gemini with an instruction:

\begin{lstlisting}
Use only this grounded data. Do not invent missing numbers.
\end{lstlisting}

\subsection{P16: Chat That Remembers}

\subsubsection*{Goal}
Support persistent named chat sessions.

\subsubsection*{Endpoints}

\begin{itemize}
  \item \texttt{GET /chat/sessions}
  \item \texttt{GET /chat/:sessionId/history}
  \item \texttt{DELETE /chat/:sessionId}
\end{itemize}

\subsubsection*{MongoDB Collections}

\textbf{sessions}

\begin{lstlisting}
sessionId
name
createdAt
lastMessageAt
\end{lstlisting}

\textbf{messages}

\begin{lstlisting}
sessionId
role: "user" | "assistant"
content
timestamp
\end{lstlisting}

\subsubsection*{Logic}

\begin{enumerate}
  \item On every chat message, load previous messages for the session.
  \item Pass recent history to Gemini.
  \item Save user message and assistant reply.
  \item Update session \texttt{lastMessageAt}.
  \item For a new session, generate a short title using the LLM.
\end{enumerate}

\subsection{P17: Frontend Dashboard}

\subsubsection*{Goal}
Build a user-facing dashboard that calls only the API gateway.

\subsubsection*{Implemented Pages}

\begin{itemize}
  \item \texttt{/login}: user login form.
  \item \texttt{/register}: registration form.
  \item \texttt{/products}: product list with pagination and category filter.
  \item \texttt{/availability}: product availability checker.
  \item \texttt{/chat}: AI assistant with session sidebar.
  \item \texttt{/profile}: user discount lookup.
  \item \texttt{/insights}: surge and peak-window analytics.
  \item \texttt{/trending}: seasonal recommendations.
\end{itemize}

\subsubsection*{Frontend API Rule}

The frontend uses:

\begin{lstlisting}
const API_BASE = "http://localhost:8000"
\end{lstlisting}

All requests go to the gateway.

\subsection{P18: Trending Today Widget}

\subsubsection*{Goal}
Show today's seasonal product recommendations.

\subsubsection*{Endpoint Used}

\begin{lstlisting}
GET /analytics/recommendations?date=TODAY&limit=6
\end{lstlisting}

\subsubsection*{Logic}

\begin{enumerate}
  \item Compute today's date dynamically in JavaScript.
  \item Fetch recommendations through gateway.
  \item Show skeleton cards while loading.
  \item Show friendly error message on failure.
  \item Render product cards with name, category, and score.
  \item Provide refresh button.
\end{enumerate}

\subsection{P19: Lean Images}

\subsubsection*{Goal}
Keep Docker final image sizes under target limits.

\subsubsection*{Optimization Logic}

\begin{itemize}
  \item Use multistage Dockerfiles.
  \item Use \texttt{npm ci --omit=dev}.
  \item Move \texttt{nodemon} to dev dependencies.
  \item Remove unused \texttt{bun} dependency from production services.
  \item Use Alpine runtime with Node.js.
  \item Add \texttt{.dockerignore} files.
  \item Copy only required runtime files.
\end{itemize}

\subsubsection*{Verified Image Sizes}

\begin{center}
\begin{tabular}{|p{5cm}|p{3cm}|p{3cm}|}
\hline
\textbf{Service} & \textbf{Target} & \textbf{Verified Size} \\
\hline
api-gateway & $<$ 150 MB & 100 MB \\
\hline
user-service & $<$ 150 MB & 118 MB \\
\hline
rental-service & $<$ 150 MB & 95.7 MB \\
\hline
analytics-service & $<$ 150 MB & 98 MB \\
\hline
agentic-service & $<$ 300 MB & 226 MB \\
\hline
frontend & $<$ 50 MB & 6.41 MB \\
\hline
\end{tabular}
\end{center}

\section{Bonus Problems}

\subsection{B1: gRPC Internal Communication}

The README describes a possible bonus where the agentic service can communicate with analytics-service or rental-service using gRPC. This project keeps HTTP communication for internal grounding calls. Therefore, this bonus is not implemented.

\subsection{B2: Graceful Rate Limit Handling}

The README describes exponential backoff with jitter for central API \texttt{429} responses. Some services handle rate limits by returning clear \texttt{429} messages or throttling central requests, but full B2 exponential backoff with jitter is not implemented everywhere.

\section{Important Algorithms Summary}

\begin{longtable}{|p{4cm}|p{5cm}|p{5cm}|}
\hline
\textbf{Problem} & \textbf{Algorithm} & \textbf{Reason} \\
\hline
P7 Availability & Sort and merge intervals & Correctly handles overlapping rentals. \\
\hline
P8 Kth busiest date & Min-heap of size K & Finds kth largest without sorting all dates. \\
\hline
P10 Free streak & Merge busy intervals and scan gaps & Finds longest free continuous date range. \\
\hline
P11 Peak window & Sliding window & Computes best 7-day period in linear time. \\
\hline
P12 Merged feed & Two-pointer sorted-list merge & Merges sorted rental streams efficiently. \\
\hline
P13 Surge days & Monotonic stack & Finds next greater rental count in linear time. \\
\hline
P14 Recommendations & Frequency map and bucket selection & Finds top seasonal products efficiently. \\
\hline
\end{longtable}

\section{How to Run}

\subsection{Local Docker}

\begin{lstlisting}
cp .env.example .env
# Fill in CENTRAL_API_TOKEN, JWT_SECRET, GEMINI_API_KEY
docker compose up --build
\end{lstlisting}

\subsection{Service URLs}

\begin{itemize}
  \item Frontend: \url{http://localhost:3000}
  \item API Gateway: \url{http://localhost:8000}
  \item User Service: \url{http://localhost:8001}
  \item Rental Service: \url{http://localhost:8002}
  \item Analytics Service: \url{http://localhost:8003}
  \item Agentic Service: \url{http://localhost:8004}
\end{itemize}

\section{Verification Checklist}

\begin{itemize}
  \item \texttt{npm run build} passes in frontend.
  \item Backend JavaScript syntax checks pass.
  \item \texttt{docker compose build} passes.
  \item \texttt{docker compose up -d} starts all services.
  \item All services become healthy.
  \item Gateway \texttt{/status} returns all downstream services as \texttt{OK}.
  \item Frontend returns HTTP 200.
\end{itemize}

\section{Conclusion}

This RentPi project is a complete microservice system for a rental platform. The API gateway provides a clean entry point, user-service manages authentication and discounts, rental-service implements product and rental algorithms, analytics-service computes demand insights, agentic-service provides grounded AI chat, and the frontend presents the system through a usable dashboard.

The most important engineering ideas are:

\begin{itemize}
  \item Protect external API tokens by routing through backend services.
  \item Use efficient algorithms for rental data problems.
  \item Cache stable data such as categories.
  \item Keep AI responses grounded in fetched data.
  \item Keep Docker images lean for reliable deployment.
\end{itemize}

\end{document}
